\chapter{Conclusions and Future Work}
\label{chapter:conclusions_and_future_work}


\section{Conclusions}




\section{Future Work}

\subsection{Integration of Meteorological Radar Data for Enhanced Rainfall Spatialization}

While the current data pipeline successfully spatializes point-based rain gauge measurements using elevation-based 
drift (e.g., Kriging with External Drift), this methodology presents inherent limitations when modeling extreme weather 
events. Elevation acts as a static covariate that effectively represents long-term, climatological precipitation trends 
(the orographic effect). However, it cannot capture the highly dynamic, localized, and heterogeneous nature of real-time 
convective storm cells. Consequently, the interpolated rainfall tensor may misrepresent the exact geographical 
distribution of the precipitation nucleus during a specific simulation step.

To achieve professional-grade meteorological realism and improve the predictive accuracy of the hydraulic simulation, 
a primary line of future work is the integration of Meteorological Radar Data (such as the C-band and S-band radar 
networks operated by the Spanish Meteorological Agency, AEMET).

This integration would fundamentally upgrade the Rainfall Engine through the following technical implementations:

\begin{itemize}
	\item \textbf{Radar Reflectivity Conversion}: Radar systems provide high-resolution spatial matrices of atmospheric 
	reflectivity ($Z$). The pipeline would need to implement algorithms to dynamically fetch these grids and convert 
	them into estimated rainfall intensity rates ($R$) using empirical formulas, such as the Marshall-Palmer relationship 
	($Z = aR^b$).
	\item \textbf{Radar-Rain Gauge Merging}: Because radar data is susceptible to systemic anomalies (e.g., beam blockage, 
	signal attenuation, or overestimation due to hail), it cannot be used in isolation. The future pipeline should 
	implement advanced merging techniques—such as Conditional Merging (CM) or Regression Kriging. In this upgraded 
	architecture, the radar imagery would serve as the highly detailed spatial covariate, while the physical rain gauges 
	would provide the "ground-truth" anchor points to correct the radar bias.
	\item \textbf{Enhanced Tensor Fidelity}: By replacing a static topographic covariate with a dynamic meteorological 
	one, the resulting $Time \times Y \times X$ data cube would precisely reflect the trajectory, shape, and intensity 
	of storm events.
\end{itemize}

Implementing this feature would drastically reduce spatial interpolation errors and provide the Cellular Automata 
($m:n-CA^k$) core with a highly realistic input tensor, which is strictly necessary for the accurate modeling of 
rapid-onset flash floods.


\subsection{GPU Acceleration via CUDA for High-Performance Hydraulic Simulation}

While the current execution framework successfully deploys the TensorFlow models via the ONNX Runtime C++ API on the 
CPU, it introduces significant computational bottlenecks when scaling to high-resolution watersheds or executing large-scale 
simulations. Flood forecasting and flash-flood modeling are inherently time-critical operations; however, the current 
CPU-bound architecture severely limits the execution throughput of the Cellular Automata ($m:n-CA^k$) core, especially when 
processing the high-dimensional spatial-temporal tensors generated by the Rainfall Engine. To transition this framework into a 
viable real-time Early Warning System (EWS), a primary future objective is the implementation of hardware-level parallelization 
leveraging NVIDIA CUDA-based GPU acceleration.

This architectural upgrade will focus on mitigating computational latency through the following technical implementations:

\begin{itemize}
	\item \textbf{ONNX Runtime CUDA Execution Provider Integration}: Currently, the exported ONNX models are evaluated using 
	the default CPU inference session. Upgrading the C++ pipeline to leverage the CUDA Execution Provider (or NVIDIA TensorRT) 
	will allow the underlying layers to execute entirely on the GPU. This will drastically reduce inference latency by exploiting 
	tensor cores for matrix operations.
	
	\item \textbf{Custom CUDA Kernels for Cellular Automata Core}: Beyond inference, the structural operations and state transition 
	rules of the Cellular Automata ($m:n-CA^k$) are highly parallelizable by nature. Because cell state updates depend primarily on 
	localized neighborhoods, developing custom CUDA kernels will allow for the simultaneous, massively parallel computation of millions 
	of grid cells, transforming the current sequential or multi-threaded CPU loops into an optimized grid-block GPU execution model.
	
	\item \textbf{Asynchronous Memory Transfer Optimization}: A critical challenge in GPU computing is the latency overhead induced 
	by Host-to-Device (CPU-to-GPU) data transfers. To prevent the PCIe bus from becoming a bottleneck, the upgraded architecture 
	will implement advanced memory management strategies, such as asynchronous data transfers via CUDA streams, the utilization of 
	pinned memory (Page-Locked Host Memory), and the strategic deployment of Unified Memory where applicable. This ensures that 
	incoming rainfall tensors and hydraulic state variables flow seamlessly between components.
\end{itemize}

Successfully shifting the computational workload from the CPU to the GPU will unlock the performance thresholds required for 
sub-real-time flood routing. This optimization will not only allow for faster-than-real-time deterministic simulations but will 
also enable the execution of concurrent, multi-scenario ensemble forecasting during critical, rapid-onset weather events.


\subsection{Multi-Resolution Spatial Aggregation and Dynamic Grid Coarsening}

The current simulation framework operates on a uniform spatial grid where computational complexity scales quadratically 
with the domain's linear dimensions; for instance, a standard $1000 \times 1000$ grid requires evaluating one million 
cells per simulation step. Although the existing pipeline implements an optimization layer that restricts inference strictly 
to hydraulically active cells (i.e., cells experiencing active precipitation or dynamic surface runoff), the underlying 
spatial domain remains locked at its maximum resolution. To further alleviate computational overhead and enable efficient 
scaling across extensive geographical basins, a critical line of future work is the implementation of a configurable, 
multi-resolution spatial aggregation framework.

This methodology will allow the simulation engine to dynamically group clusters of adjacent cells into macro-cells (e.g., 
$2\times2$, $3\times3$, or $4\times4$ configurations) according to user-defined scalability requirements. The upgrade will 
be structured around the following core technical developments:

\begin{itemize}
	\item \textbf{Dynamic Grid Resolution Scaling}: The engine will support arbitrary spatial coarsening factors ($k \times k$). 
	By aggregating the fine-grained topographic and hydraulic tensors into coarser representations, the total number of potential 
	computational nodes is reduced by a factor of $k^2$, drastically decreasing both the memory footprint and the number of 
	state-transition evaluations per timestep.
	
	\item \textbf{Synergistic Active-Cell Masking Integration}: The proposed aggregation layer will natively interface with the 
	existing active-cell tracking mechanism. Instead of evaluating individual active pixels, the system will dynamically construct 
	a macro-cell activity mask. A macro-cell will be flagged as active and queued for computation only if at least one of its 
	underlying high-resolution constituent cells exhibits hydraulic activity, ensuring that computational resources are exclusively 
	allocated to evolving flood fronts.
	
	\item \textbf{Mass-Conservative Mapping Operators}: Transitioning between spatial scales introduces the risk of numerical 
	diffusion and mass-balance violations. Future implementations must develop mathematically rigorous mapping operators. 
	Downscaling (aggregation) must employ mass-conservative averaging for topographic elevation and input rainfall volumes, while 
	upscaling (disaggregation) must utilize specialized interpolation techniques to accurately reconstruct high-resolution flood 
	boundaries from macro-cell hydraulic heads without artificial water loss or creation.
\end{itemize}

By integrating dynamic spatial aggregation with the existing active-domain tracking, the framework will achieve unprecedented execution 
speeds. This will allow users to seamlessly trade off a controllable margin of spatial resolution in exchange for the ultra-high 
throughput necessary for real-time ensemble flash-flood forecasting.
